\section{Summable parameter estimates for the sector expansion}
\label{sec:sector-analytic-prices}
We now combine the real-parameter interpolation with the local
Gaussian and connected-support estimates. The important point is
that one uses the same supports, labels and protected fine reads
throughout the real parameter family. A separate estimate at each
parameter, with an unrelated choice of labels, would not bound a
finite linear combination of the resulting activities.

Fix the group, the block side $L$, a derivative order $j$, a
nonnegative integer $m$ measuring the spatial weight in the derivative
budget, and all the positive safety and support reserves. Write
$y=\beta_f(x)$ for the reference fine coupling, which tends to
infinity with the scale parameter $x$. A nonnegative function has
upper exponent $a$ below if, for each $a'>a$, it is eventually at
most $y^{a'}$. Finite sums, products and fixed powers have the
corresponding upper-exponent bounds. Constants, local cover counts,
and all the fixed-order derivative factors of upper exponent zero
may therefore be absorbed using a strict exponent reserve.

The incoming sector is assumed to satisfy the local hypotheses of
Sections~\ref{sec:activity-jet-assembly} and
\ref{sec:gaussian-differentiation}, uniformly for all real points
$|t|\le6R$ in the coupling band. Uniformity is meant in the
following sense. The quantitative local data of those sections (the
local jets of the four species, the safety inequality, the
compensated chart and Jacobian bounds with their derivative scales,
and the smoothness of the local Gaussian density) are given with one
set of constants for the whole family, and not by a separate choice
at each $t$. The order of choices is: first the lower bound on the
block side; then the block side, the stage orders, the positive
reserves and the other fixed data listed above; then one threshold
in $x$, valid for all allowed periods, sectors and real parameter
points. The path, clipping rule, ordinary
and forced support dictionaries, and held fine configurations are
common. The coupling is $\beta(t)=\beta_0\ell(t)^2$, where
\eqref{eq:real-affine-band} holds. The incoming ordinary activity
values and forced activity values on each fixed support are
polynomials of degree at most two on $[-1,1]$. We take
\begin{equation}\label{eq:sector-radius-choice}
                       R=y^{1/256},\qquad \delta=y^{-1/128}.
\end{equation}
The geometric radius $B(x)$ and its ball count $b_B=(2B+1)^4$
have upper exponent zero. Fix the incoming support exponent $E$ and
the desired outgoing exponent $e$ with $E>e\ge0$.
All estimates below are eventually uniform in the allowed periods,
sectors and configurations. The threshold is chosen after the fixed
data just listed. Construction of those incoming sectors and of
their uniform safety estimates is still required; the purpose of
this section is to prove the analytic output estimates from them.

\subsection{Common prices before integration}
The original optional-activity rooted allowance has upper exponent
$a\le-7/16$. To obtain one price valid along the parameter family,
sum the nonnegative prices at the three nodes $-1,0,1$ for each
ordinary and forced label. The sum over all such labels through a
root increases by at most a factor three, and each individual
price is at most three times the old rooted allowance. This last
assertion uses nonempty supports and nonnegative support weights.

The complete local logarithmic factor is still polynomial of degree
at most two: the old vacuum and action terms either acquire the factor
$\ell(t)^2$ or remain constant, the incoming regular occurrence is
quadratic by hypothesis, and the enforcement factor uses held fine
reads and is parameter-independent. The forced local factor likewise
is its incoming quadratic value multiplied by a fixed enforcement factor.
Apply the small-activity estimate of
Section~\ref{sec:real-coupling-interpolation} to the ordinary
exponential-minus-one factors. Its smallness conditions are satisfied
eventually: if $g(x)$ is the incoming allowance, then
$3g(x)\to0$ and $R^2g(x)$ has negative upper exponent. For a forced factor use
ordinary degree-two interpolation at the same three nodes; there
is no logarithm to take. The resulting common prices cost at most
a fixed constant times $1+R^2$. Their upper exponent is at most
\[
                 a+\frac{2}{256}
                     \le-\frac{55}{128}<-\frac25.
\]
Thus both are paid by an enlarged target with upper exponent
$-2/5$, after increasing the scale threshold. This explains the
specific reserve in \eqref{eq:sector-radius-choice}.

For seed packets, use the larger-real-node estimate with the full
stock weight. If a packet contains $s\ge1$ seeds, its real price
has the form
\[
                    \exp\{(C\mathfrak d^5-\beta\mu\tau^2)s\},
\]
where $\mathfrak d\ge1$ is the local geometric control function.
Its stock-coordinate count satisfies $n\le C_1s\mathfrak d^5$.
The interpolation factor is $K_j(1+\delta)^n$.
Since $0\le\delta\le1$, $(1+\delta)^n\le e^n$, and
$K_j\le e^{K_j}$ for $K_j\ge0$. Hence increasing $C$ to
$C+C_1+K_j$ absorbs the entire factor, because
$s\mathfrak d^5\ge1$. The negative packet energy is unchanged.
The absorption and root-count argument of
Section~\ref{sec:packet-estimates} consequently still applies.

Warnings are constant in this parameter when evaluated with the
specified held fine reads and rescaled fluctuation coordinate.
Their coarse-field jets retain the common reduced-margin excursion
indicator of Section~\ref{sec:activity-jet-assembly}. No derivative
of that indicator is taken. Products of the four species thus have
one common nonnegative coefficient majorant, with the same
excursion event and the same support set as before interpolation.

There is one further factor from the local Gaussian ratio and the
local stock coordinate count. Let $I(S)$ be the set of fine bonds
read by the local term with coarse support $S$, and let $d_G$ be
the real color dimension. Its combined power is bounded by
$2|I(S)|d_G$. If $c_B(S)$ is the coarse cover number, then
\[
 |I(S)|\le |S|\,\operatorname{dens}(L)
       \le c_B(S)b_B\operatorname{dens}(L),
                         \qquad b_B=(2B+1)^4.
\]
For any fixed $\alpha>0$, \eqref{eq:sector-radius-choice} gives,
eventually,
$2\delta b_B\operatorname{dens}(L)d_G\le\alpha$.
Using $1+\delta\le e^\delta$ proves
\begin{equation}\label{eq:sector-local-dimension}
                 (1+\delta)^{2|I(S)|d_G}
                                   \le e^{\alpha c_B(S)}.
\end{equation}
Only the dimension of the local read set has been charged.
The factor is paid by a support-exponent reserve of size $\alpha$.

\subsection{The moment estimate after a finite parameter test}
Let $A$ be a connected finite set of primitive labels. Denote its
assigned union support by $S_A$ and its warning set by $W_A$. Its
shape charge is the total cost of
Section~\ref{sec:mixed-activity-prices},
\[
             s_A=\sum_{l\in A}\bigl(c_B(S_l)+d_l\bigr),
\]
where $S_l$ is the support of the label $l$ and $d_l\ge0$ its extra
cost: the purchased fixed multiple of its number of seeds for a
packet, and the purchased optional cost for an optional label. Let
$a_l\ge0$ be the common
primitive price for label $l$, with its warning gain included only
after Gaussian integration. Put $P_A=1+s_A\ge1$. The estimates
just proved, the compensated Jacobian estimates, and the product
rule give a coefficient bound for every derivative through order
$j$ of the local integrand. After a finite test $\lambda$, its
norm is bounded by
\begin{equation}\label{eq:sector-integrand-test}
 M_R(\lambda)\,C_j e^{\alpha c_B(S_A)}
       \left(\prod_{l\in A}a_l^{\circ}\right)
       P_A^j\,\mathbf1_{\mathcal E_A}(z)\,F_{I(A)}(z)^{Kj}.
\end{equation}
Here $a_l^{\circ}$ is the price before inserting the warning
probability gain, $\mathcal E_A$ is the joint reduced-margin
excursion event, $F_{I(A)}\ge1$ is the Gaussian coordinate
control function on the local read set, and $K$ is the fixed
Jacobian polynomial order. The constant $C_j$ in this display
may be replaced by a nonnegative function of upper exponent zero;
this includes the compensated derivative scales. More explicitly,
if those scales are $r$ and $L'F_{I(A)}^K$, their product-rule
cost is bounded by
\[
 \{1+(2^KL'+2r)\}^{j}P_A^jF_{I(A)}^{Kj}.
\]
Indeed $F_{I(A)},P_A\ge1$ and $s_A\le P_A$, so
$1+2^KL'F_{I(A)}^K+2rs_A$ is at most
$\{1+(2^KL'+2r)\}P_AF_{I(A)}^K$. Fixed polarization constants
are included in $C_j$.

For clarity, the probabilistic step in this estimate is as follows.
The local Gaussian moment calculation in
Section~\ref{sec:gaussian-moment-prices} supplies an integer $k$
and a function $V(x)$ of upper exponent zero such that
\[
 \left(\int F_{I(A)}^{2Kj}\,d\nu\right)^{1/2}
                                  \le(1+|V(x)|)P_A^k.
\]
The common warning estimate is
$\nu(\mathcal E_A)\le\rho^{2|W_A|}$.
Cauchy--Schwarz therefore gives
\[
 \int\mathbf1_{\mathcal E_A}F_{I(A)}^{Kj}\,d\nu
                  \le(1+|V(x)|)P_A^k\rho^{|W_A|}.
\]
This is the source of the warning price in $a_l$; no independence
of the warning events is assumed. Combining with
\eqref{eq:sector-integrand-test}, and then averaging the finitely
many forest and tagging choices, gives the derivative price
\begin{equation}\label{eq:sector-polymer-price}
 p_A=\pi(x)e^{\alpha c_B(S_A)}
                   \left(\prod_{l\in A}a_l\right)P_A^{k+j},
\end{equation}
where $\pi$ has upper exponent zero and is independent of
$\lambda$. The bound on the tested derivative is
$M_R(\lambda)p_A$.

Passing derivatives through the integral is justified by the
nearby integrable bounds in Section~\ref{sec:gaussian-differentiation}.
There are only finitely many parameter points in a test, so
finite summation commutes with derivatives and integrals. A finite
average over forests preserves the same bound. On healthy supports
the protected gauge average equals the unaveraged local expression
as a germ: the safety plateau is open and the gauge change leaves
the protected expression invariant. Thus the estimate is for the
actual coarse activity, including its derivatives. An inactive
support gives the zero germ. All of these comparisons use the common
support dictionary, not a parameter-dependent relabeling.

The mixed connected-set estimate \eqref{eq:mixed-reserved-tail}
of Section~\ref{sec:mixed-activity-prices}, applied to the enlarged
primitive prices, gives
\begin{equation}\label{eq:sector-first-root}
       \sup_z\sum_{A:z\in S_A}p_Ae^{(E-\alpha)c_B(S_A)}
                                              \le\pi(x)u_{k+j}(x).
\end{equation}
Here $E$ is the incoming support exponent. For an integer $n\ge0$,
$u_n(x)$ denotes the connected allowance of the enlarged target with
polynomial order $n$: the right side of \eqref{eq:mixed-ceiling} for
that target, with no restriction on the label count, the total cost
or the number of seeds. The vertex condition of that estimate is the
purchase of the enlarged target made in the first subsection; the
factor $P_A^{k+j}=(1+s_A)^{k+j}$ in $p_A$ is its polynomial in the
total cost. The allowance $u_{k+j}$ has upper exponent $-2/5$. To see
the support bookkeeping,
the left side contains
$e^{\alpha c_B(S_A)}e^{(E-\alpha)c_B(S_A)}
=e^{Ec_B(S_A)}$. Thus the mixed connected-set estimate is applied
with its original support weight, and its shape moment is exactly
$k+j$. This invokes the same connected-set and moment bound used
for the non-parameter sector, with the enlarged primitive prices
just constructed.

\subsection{The second connected expansion}
Recover an absolute parameter series for each derivative from the
finite-test estimate with radius $R$. Section~\ref{sec:background-ray-estimates}
gives radius $R/2$ and price $2p_A$, once $R\ge2$.
The argument can first be applied to each scalar derivative;
the fixed-order polarization bound controls the full derivative.
The derivatives of the outgoing budget are iterated group increments
in directions of norm at most the threshold $\theta$, whereas the
local estimates are stated in the compensated charts. The passage
between the two is part of the uniform local hypotheses: appending
one such increment increases the controlled chart scale by at most a
fixed multiple $c_v$ of the threshold scale, and the chart scale $v$
is chosen with $jc_v\le v$, as in
Section~\ref{sec:parameter-comparison-suppliers}. Thus the at most
$j$ increments of a derivative of the budget stay inside the chart
estimates used above.
A logarithmic cluster with $q$ occurrences is a product of these
primitive integrated activities times its Ursell coefficient
and $1/q!$. Derivatives through order $j$ cost at most $q^j$.
Since $q\le2^q$ for $q\ge1$, this is paid by multiplying each
vertex price by $2^j$. Consequently the absolute cluster price is
\[
       w_\Gamma=\frac{|u(\Gamma)|}{q!}
                           \prod_{r=1}^{q}(2^{j+1}p_{A_r}).
\]
This includes repeated labels and their tuple factorial exactly
as in Section~\ref{sec:cluster-expansion}.

Fix the desired outgoing exponent $e$ with $0\le e<E$, and set
$\alpha=(E-e)/4>0$. Put
\[
                   Y(x)=(1+2^{j+1}\pi(x))u_{k+j}(x).
\]
Equation~\eqref{eq:sector-first-root} pays the enlarged vertex
prices at exponent $e+2\alpha$, because
$e+2\alpha\le E-\alpha$. The connected-tree degree criterion is
\[
                         \frac{4b_BY(x)}{\alpha}\le\frac12.
\]
It holds eventually: $b_B$ has upper exponent zero and $Y$ has
upper exponent $-2/5$. The cluster estimate then gives
\[
                \sup_z\sum_{\Gamma:z\in S_\Gamma}
                      w_\Gamma e^{(e+\alpha)c_B(S_\Gamma)}
                                                       \le2Y(x).
\]

The spatially weighted derivative budget at a root costs at most
$\{s_B(1+c_B(S_\Gamma))\}^{mj}$, where
$s_B=4(2B+1)$. To verify the geometry, connect the balls of a
minimal connected cover. A path through at most $c_B(S_\Gamma)$
balls bounds the distance between two sites of $S_\Gamma$ by
$c_B(S_\Gamma)(4B+2)+2B$, which is at most
$s_B(1+c_B(S_\Gamma))$. Locality allows the field variation to
be restricted to this support. The weighted direction norm on
that support is therefore increased by the stated distance
factor, and taking at most $j$ derivatives gives the power $mj$.
For every fixed integer $d\ge0$,
\[
       P_{\alpha,d}:=\sup_{r\ge0}(1+r)^d e^{-\alpha r}<\infty.
\]
For example, expand $(1+r)^d$ and bound each $r^l$ by
$l!\alpha^{-l}e^{\alpha r}$, using the exponential series.
The extra spatial factor is thus paid by the remaining
$\alpha$ reserve. Including the fixed polarization factor $C_j'$
we obtain the final rooted derivative price
\begin{equation}\label{eq:sector-ordinary-final}
                       2C_j's_B^{mj}P_{\alpha,mj}Y(x).
\end{equation}
It has upper exponent $-2/5$. Since $-2/5<-3/8$, it is eventually
at most $y^{-3/8}$. This is the ordinary parameter-test allowance
at radius $R/2$, on the actual connected-cluster dictionary.

\subsection{The forced common majorant}
For forced terms no coarse derivatives are required. The same
finite-test and radius-halving argument gives ordinary constituent
series prices with rooted allowance $2u_0(x)$ and forced
constituent prices with rooted allowance
$2u_0(x)e^{-M(x)}$. The nonnegative mark $M(x)$ is the seed mark of
the enlarged target: the mark per seed which the mixed estimate
\eqref{eq:mixed-ceiling} extracts once for every seed of a
component, and which is paid from the seed expense of the packet
prices. It is available here because every forced
constituent has a nonempty flag inventory, each flag supplies a
seed in one of its packets, and hence its components contain at
least one seed. It is distinct from the warning probability gain
$\rho^{|W_A|}$, which has already been included in the primitive
prices $a_l$ and is present in ordinary and forced constituents
alike. Form the exclusion-correction activities exactly as
in Section~\ref{sec:forced-exclusion}. Each is a finite product
of forced factors and powers of regular cluster factors, with the
same multiplicity factorials. Absolute coefficient norms multiply,
so \eqref{eq:exclusion-output-price} is also a majorant for the
whole parameter series at radius $R/2$.

The remaining eligibility and flag indicators depend on the fixed
coarse configuration and common supports, and are independent of
the parameter. Multiplication by such an indicator only decreases
the coefficient bound. This yields one majorant per final support,
valid for all coarse configurations, rather than a separate
majorant selected after evaluating a configuration.

Apply the forced-exclusion bound with its $u=2u_0(x)$ and the
reserve $\alpha$. Its second-expansion allowance is
$Z(x)=20u=40u_0(x)$. Both $Z\le1/2$ and
$4b_BZ/\alpha\le1/2$ hold eventually. Hence the final common
majorants $k_D$ satisfy
\begin{equation}\label{eq:sector-forced-final}
        \sup_z\sum_{D:z\in S_D}k_De^{e c_B(S_D)}
                            \le e^{-M(x)}\frac{Z(x)}{1-1/2}.
\end{equation}
This keeps the compulsory exponential seed gain through arbitrarily
many regular exclusion corrections. In particular, the parameter
extension has not consumed that gain.

Equations~\eqref{eq:sector-ordinary-final} and
\eqref{eq:sector-forced-final} are the broad analytic estimates
for the sector output before its finite Taylor and local
projection operations. They use the uniform local sector estimates,
the same connected-support sums, and the same Gaussian warning
bounds as the real one-step construction. The later projection
prices and the construction of the allowed coupling family are
separate inputs when these estimates are applied to the selected
background history.
